\section*{Overview}

Minimum spanning tree is the weighted-graph tool for ``connect everything as cheaply as possible.'' In contest code,
Kruskal plus DSU is usually the shortest reliable route, and it also exposes the structural properties that later MST
reductions depend on.

\section*{When to Use It}

Use MST ideas when:

\begin{itemize}
  \item the graph is undirected and weighted,
  \item the goal is to connect all vertices with minimum total cost,
  \item the statement hides a connectivity threshold or bottleneck argument.
\end{itemize}

\section*{Core Idea}

A spanning tree uses exactly \(n-1\) edges and connects every vertex. Kruskal's algorithm sorts edges by weight and
adds an edge exactly when it joins two different DSU components.

\section*{Key Insight}

The reason Kruskal is safe is the cut property: for any cut of the graph, the lightest edge crossing that cut belongs to
some MST. Kruskal keeps applying this fact globally by always taking the cheapest edge that does not create a cycle.

\section*{Worked Problem}

\paragraph{Problem.}
There are \(n\) outposts and weighted undirected roads. Build roads so every outpost can reach every other and the total
construction cost is minimum.

\paragraph{Algorithm outline.}

\begin{itemize}
  \item Sort all roads by weight.
  \item Scan from cheapest to most expensive.
  \item If a road joins two different DSU components, take it.
  \item Stop after taking \(n-1\) roads.
\end{itemize}

\paragraph{Why this greedy step is correct.}
At the moment Kruskal considers edge \(e\), its endpoints lie in two different components of the partial forest. Those
components define a cut, and \(e\) is the lightest edge crossing that cut that has not already been skipped for creating
cycles. By the cut property, some MST can include it.

\section*{Example Problem Pattern}

Another common use is a threshold problem:

\begin{quote}
What is the minimum value \(T\) such that the graph becomes connected using only edges of weight at most \(T\)?
\end{quote}

The answer is the maximum edge weight used by the MST. This is a bottleneck interpretation of the same structure.

\section*{Correctness Intuition}

Every time Kruskal adds an edge, it is forcing one safe edge into the final tree. Repeating that exchange argument
eventually yields a full spanning tree, and no cheaper tree can exist because each chosen edge was safe at the moment it
was taken.

\section*{Complexity Analysis}

\begin{itemize}
  \item sorting edges: \(O(m \log m)\),
  \item DSU operations: near-linear,
  \item total: \(O(m \log m)\).
\end{itemize}

\section*{Implementation}

The reference code returns both the total MST weight and the chosen edge list. That makes it easy to continue with tree
queries or bottleneck analysis afterward.

\section*{Common Pitfalls}

\begin{itemize}
  \item Using MST language on a directed graph.
  \item Forgetting to detect the disconnected case.
  \item Confusing MST with single-source shortest paths.
  \item Using 32-bit sums when edge weights are large.
\end{itemize}

\section*{Variants / Extensions}

\begin{itemize}
  \item Prim's algorithm for dense graphs or adjacency-matrix settings.
  \item Maximum spanning tree.
  \item Second-best MST and replacement-edge queries.
\end{itemize}

\section*{Practice Problems}

\begin{itemize}
  \item Standard minimum spanning tree construction.
  \item Threshold-to-connect problems.
  \item Build the MST first, then answer queries on the resulting tree.
\end{itemize}

\section*{References}

\begin{itemize}
  \item \href{https://cp-algorithms.com/graph/mst_kruskal.html}{cp-algorithms: Minimum Spanning Tree - Kruskal}.
  \item \href{https://usaco.guide/CPH.pdf}{Competitive Programmer's Handbook}.
  \item \href{https://codeforces.com/catalog}{Codeforces Catalog}.
\end{itemize}
